\chapter*{Kurzfassung} %*-Variante sorgt dafür, das Abstract nicht im Inhaltsverzeichnis auftaucht
\phantomsection\pdfbookmark{Kurzfassung}{kurzfassung}

\chapter*{Kurzfassung}

Im Rahmen dieser Studienarbeit wurde ein bereits bestehender Demonstrationsroboter für den Einsatz auf Hochschulmessen weiterentwickelt. Ziel war es, den im Rahmen einer vorherigen Studienarbeit entstandenen Funktionsprototypen sowohl technisch als auch gestalterisch weiter auszubauen und dessen Interaktionsmöglichkeiten mit Messebesuchern zu verbessern.

Hierfür wurden die mechanische Konstruktion, das Gehäusedesign sowie die Mensch-Roboter-Interaktion überarbeitet. Zu den wesentlichen Weiterentwicklungen zählen die Neugestaltung des Fahrwerks zur Erhöhung der Standsicherheit, die Entwicklung eines modernen Gehäuses mit humanoiden Gestaltungselementen, die Integration beweglicher Roboterarme sowie einer Hinderniserkennung auf Basis von Ultraschallsensoren. Darüber hinaus wurde ein KI-gestütztes Sprachdialogsystem implementiert, das eine Kommunikation zwischen Roboter und Besuchern ermöglicht.

Die konstruierten Komponenten wurden überwiegend mittels additiver Fertigung hergestellt und anschließend in das bestehende Robotersystem integriert.

Die Validierung erfolgte im Rahmen zweier Messeeinsätze unter realen Einsatzbedingungen. Dabei zeigte sich, dass insbesondere das äußere Erscheinungsbild und die Interaktionsmöglichkeiten des Roboters von den Besuchern sehr positiv aufgenommen wurden. Gleichzeitig konnten technische Schwachstellen hinsichtlich der Hardwarezuverlässigkeit, der Energieversorgung sowie der Sprachinteraktion identifiziert werden. Die gewonnenen Erkenntnisse bilden die Grundlage für die weitere Entwicklung des Roboters und zukünftige Studienarbeiten.

\clearpage